%% ============================================================
%%  Rapport de Module : Big Data — Architecture CDC Medallion
%%  Auteur : Zakaria Elghouali
%%  Langue : Français
%% ============================================================
\documentclass[12pt,a4paper]{report}

%% ── Encodage & Langue ──────────────────────────────────────
\usepackage[utf8]{inputenc}
\usepackage[T1]{fontenc}
\usepackage[french]{babel}

%% ── Mise en page ───────────────────────────────────────────
\usepackage[top=2.5cm, bottom=2.5cm, left=3cm, right=2.5cm]{geometry}
\usepackage{setspace}
\onehalfspacing
\setlength{\parindent}{0pt}
\setlength{\parskip}{6pt}

%% ── Polices ────────────────────────────────────────────────
\usepackage{lmodern}
\usepackage{microtype}

%% ── Couleurs & Graphiques ──────────────────────────────────
\usepackage{xcolor}
\usepackage{graphicx}
\usepackage{tikz}
\usetikzlibrary{shapes.geometric, arrows.meta, positioning, fit, backgrounds, calc}

%% ── Tables ─────────────────────────────────────────────────
\usepackage{booktabs}
\usepackage{tabularx}
\usepackage{longtable}
\usepackage{multirow}
\usepackage{array}

%% ── Code ────────────────────────────────────────────────────
\usepackage{listings}

\lstset{
    basicstyle=\ttfamily\small,
    backgroundcolor=\color{gray!10},
    frame=single,
    rulecolor=\color{gray!40},
    breaklines=true,
    breakatwhitespace=true,
    keywordstyle=\color{blue!70!black}\bfseries,
    stringstyle=\color{green!50!black},
    commentstyle=\color{gray!70}\itshape,
    numbers=left,
    numberstyle=\tiny\color{gray},
    numbersep=6pt,
    tabsize=2,
    showstringspaces=false,
    captionpos=b,
}

%% ── Liens & PDF ─────────────────────────────────────────────
\usepackage{hyperref}
\hypersetup{
    colorlinks=true,
    linkcolor=blue!60!black,
    urlcolor=blue!70!black,
    citecolor=green!50!black,
    pdftitle={Rapport Big Data — Architecture CDC Medallion},
    pdfauthor={Zakaria Elghouali},
}

%% ── En-têtes ────────────────────────────────────────────────
\usepackage{fancyhdr}
\pagestyle{fancy}
\fancyhf{}
\fancyhead[L]{\small\textit{Architecture CDC Medallion}}
\fancyhead[R]{\small\textit{Module Big Data}}
\fancyfoot[C]{\thepage}
\renewcommand{\headrulewidth}{0.4pt}

%% ── Boîtes ──────────────────────────────────────────────────
\usepackage{tcolorbox}
\tcbuselibrary{skins, breakable}
\newtcolorbox{infobox}[1][]{
    enhanced, breakable,
    colback=blue!5!white, colframe=blue!60!black,
    fonttitle=\bfseries, title=#1,
}

%% ── Utilitaires ─────────────────────────────────────────────
\usepackage{enumitem}
\usepackage{caption}
\usepackage{amsmath}

%% ── Couleurs personnalisées ──────────────────────────────────
\definecolor{bronze}{RGB}{205,127,50}
\definecolor{argent}{RGB}{150,150,160}
\definecolor{or}{RGB}{200,160,0}
\definecolor{darkblue}{RGB}{0,51,102}
\definecolor{kafkared}{RGB}{200,50,50}
\definecolor{sparkblue}{RGB}{30,100,200}
\definecolor{trinoviolet}{RGB}{120,80,180}

%% ════════════════════════════════════════════════════════════
\begin{document}

%% ── Page de Titre ───────────────────────────────────────────
\begin{titlepage}
    \begin{center}
        \vspace*{0.5cm}
        {\Large\bfseries École Nationale Supérieure d'Électrotechnique,\\
        d'Électronique, d'Informatique, d'Hydraulique\\
        et des Télécommunications}\\[0.3cm]
        {\large\bfseries ENSET Mohammedia}\\[0.3cm]
        \rule{\linewidth}{0.6pt}\\[0.5cm]

        {\large Module :}\\[0.2cm]
        {\LARGE\bfseries Big Data \& Traitement des Données Massives}\\[1cm]

        \rule{\linewidth}{1.5pt}\\[0.8cm]

        {\fontsize{20}{24}\selectfont\bfseries
        Pipeline CDC en Temps Réel\\[0.3cm]
        Architecture Medallion sur Lakehouse}\\[0.5cm]
        {\large\itshape Change Data Capture $\cdot$ Apache Kafka $\cdot$
        Apache Spark $\cdot$ Apache Iceberg}

        \vspace{0.8cm}
        \rule{\linewidth}{1.5pt}

        \vspace{1.5cm}
        \begin{minipage}{0.45\textwidth}
            \begin{flushleft}
                \textbf{Réalisé par :}\\[0.3cm]
                Zakaria \textsc{Elghouali}\\[0.1cm]
                \small\href{mailto:zakariaelgouali@gmail.com}
                           {zakariaelgouali@gmail.com}
            \end{flushleft}
        \end{minipage}
        \hfill
        \begin{minipage}{0.45\textwidth}
            \begin{flushright}
                \textbf{Encadrant :}\\[0.3cm]
                M. \textsc{[Nom Encadrant]}\\[0.6cm]
                \textbf{Année universitaire :}\\[0.1cm]
                2025 -- 2026
            \end{flushright}
        \end{minipage}

        \vfill
        {\small\today}
    \end{center}
\end{titlepage}

%% ── Résumé ──────────────────────────────────────────────────
\chapter*{Résumé}
\addcontentsline{toc}{chapter}{Résumé}

Ce rapport présente la conception et l'implémentation d'un pipeline de
\textit{Change Data Capture} (CDC) en temps quasi-réel, reposant sur une
\textbf{architecture Medallion} (Bronze / Argent / Or). Le projet capte les
modifications de lignes effectuées dans une base de données PostgreSQL
simulant un système d'e-commerce et les propage, via Apache Kafka et
Debezium, jusqu'à un \textit{data lakehouse} basé sur Apache Iceberg et MinIO.

Les données traversent trois couches successives traitées par Apache Spark
Structured Streaming : la couche \textbf{Bronze} (ingestion brute, immuable),
la couche \textbf{Argent} (déduplication et typage fort) et la couche
\textbf{Or} (agrégats métier). Le moteur SQL distribué Trino~435 permet
d'interroger toutes les couches en SQL standard. L'ensemble est exposé via
une API REST FastAPI sécurisée par JWT, un tableau de bord interactif React,
et supervisé par Prometheus, Grafana et Loki.

La solution est entièrement conteneurisée (Docker Compose, 21 services),
déployée sur une machine virtuelle Oracle Cloud Infrastructure (Ubuntu~22.04,
11,7~Go RAM), et conçue pour être reproductible et opérationnelle en moins de
15~minutes sur n'importe quelle machine disposant de Docker.

\bigskip
\textbf{Mots-clés :} Change Data Capture, Medallion Architecture, Apache
Kafka, Debezium, Apache Spark, Apache Iceberg, Trino, Data Lakehouse,
Big Data, Temps Réel, Docker.

%% ── Tables ──────────────────────────────────────────────────
\tableofcontents
\listoffigures
\listoftables

%% ════════════════════════════════════════════════════════════
\chapter{Introduction}
%% ════════════════════════════════════════════════════════════

\section{Contexte et motivation}

Dans les systèmes d'information modernes, les données évoluent en permanence :
commandes créées, clients mis à jour, paiements validés. Il est essentiel de
propager ces changements vers les systèmes analytiques \emph{sans interrompre
les opérations transactionnelles}. Les approches classiques (extraction en
batch toutes les nuits) introduisent une latence inacceptable pour les KPI
temps réel et peuvent surcharger la base source.

Le paradigme du \textbf{Change Data Capture} (CDC) répond à ce besoin : plutôt
que de copier l'intégralité d'une table à intervalles réguliers, le CDC capture
chaque modification --- insertion, mise à jour, suppression --- dès qu'elle
survient, avec une latence inférieure à la seconde.

Ce projet implémente un pipeline CDC complet à partir d'une base PostgreSQL
d'e-commerce et construit un \textit{lakehouse} analytique selon le modèle
Medallion, en s'appuyant exclusivement sur des technologies open-source
leaders dans l'industrie.

\section{Objectifs}

\begin{enumerate}
    \item Mettre en \oe{}uvre la capture de modifications en temps quasi-réel
          avec une latence end-to-end inférieure à 10~minutes.
    \item Construire une architecture Medallion en trois couches (Bronze /
          Argent / Or) avec Apache Iceberg comme format de table ouvert.
    \item Démontrer la déduplication, le typage fort et l'agrégation métier
          dans un contexte de flux de données continu.
    \item Exposer les données via une API REST sécurisée (JWT) et un tableau
          de bord interactif.
    \item Déployer l'ensemble sur une infrastructure conteneurisée
          auto-suffisante (Docker Compose, 21 services).
\end{enumerate}

\section{Organisation du rapport}

\begin{itemize}[leftmargin=2cm]
    \item \textbf{Chapitre~2} --- État de l'art : CDC, architecture Medallion,
          technologies clés.
    \item \textbf{Chapitre~3} --- Architecture globale du système.
    \item \textbf{Chapitre~4} --- Implémentation du pipeline de données.
    \item \textbf{Chapitre~5} --- API REST, interface utilisateur et monitoring.
    \item \textbf{Chapitre~6} --- Déploiement et résultats obtenus.
    \item \textbf{Chapitre~7} --- Conclusion et perspectives.
\end{itemize}

%% ════════════════════════════════════════════════════════════
\chapter{État de l'art}
%% ════════════════════════════════════════════════════════════

\section{Change Data Capture (CDC)}

Le CDC est une technique qui identifie et capture les changements apportés à
une base de données relationnelle, puis les transmet aux systèmes en aval. Il
existe plusieurs approches, résumées dans le tableau~\ref{tab:cdc_compare}.

\begin{table}[h]
\centering
\caption{Comparaison des approches CDC}
\label{tab:cdc_compare}
\renewcommand{\arraystretch}{1.3}
\begin{tabular}{lcccc}
\toprule
\textbf{Approche} & \textbf{Latence} & \textbf{Charge source}
                  & \textbf{Fidélité} & \textbf{Complexité} \\
\midrule
Snapshot complet   & Heures      & Très haute & Partielle & Faible  \\
Timestamp-based    & Minutes     & Moyenne    & Partielle & Faible  \\
Trigger-based      & Secondes    & Haute      & Complète  & Moyenne \\
Log-based (WAL)    & Millisecondes & Très faible & Complète & Haute  \\
\bottomrule
\end{tabular}
\end{table}

L'approche \textbf{log-based} est la plus performante car elle lit directement
le journal des transactions (WAL) sans ajouter de charge à la base source. Elle
capture toutes les opérations, y compris les suppressions, et garantit l'ordre
exact des modifications. PostgreSQL expose ce mécanisme via la réplication
logique (\texttt{wal\_level = logical}).

\section{Architecture Medallion}

Popularisée par Databricks, l'architecture Medallion organise un
\textit{data lakehouse} en trois couches de qualité croissante
(figure~\ref{fig:medallion}).

\begin{figure}[h]
\centering
\begin{tikzpicture}[node distance=2.5cm, >=Stealth]
    \tikzstyle{layer} = [rectangle, rounded corners=8pt,
        minimum width=3.4cm, minimum height=1.4cm,
        text centered, font=\bfseries\small, draw, thick]

    \node[layer, fill=bronze!25, draw=bronze!80] (B)
        {BRONZE\\[2pt]{\footnotesize Données brutes}};
    \node[layer, fill=argent!30, draw=argent!80, right of=B, xshift=0.8cm] (S)
        {ARGENT\\[2pt]{\footnotesize Données propres}};
    \node[layer, fill=or!40, draw=or!90, right of=S, xshift=0.8cm] (G)
        {OR\\[2pt]{\footnotesize Agrégats métier}};

    \draw[->, ultra thick, bronze!70] (B) -- (S)
        node[midway, above, font=\footnotesize\itshape] {MERGE INTO};
    \draw[->, ultra thick, argent!60] (S) -- (G)
        node[midway, above, font=\footnotesize\itshape] {Agrégation};

    \node[below=0.5cm of B, text width=3.4cm, font=\footnotesize, align=center]
        {Append-only\\ Audit trail\\ Latence : 30 s};
    \node[below=0.5cm of S, text width=3.4cm, font=\footnotesize, align=center]
        {Upsert ACID\\ Schéma typé\\ Latence : 60 s};
    \node[below=0.5cm of G, text width=3.4cm, font=\footnotesize, align=center]
        {KPI métier\\ Snapshot quotidien\\ Latence : 5 min};
\end{tikzpicture}
\caption{Les trois couches de l'architecture Medallion}
\label{fig:medallion}
\end{figure}

\section{Apache Iceberg --- Format de table ouvert}

Apache Iceberg est un format de table analytique open-source conçu pour
surmonter les limitations du format Hive classique~:

\begin{itemize}
    \item \textbf{Transactions ACID} : lectures et écritures cohérentes sans
          verrou global.
    \item \textbf{Évolution de schéma} : ajout ou renommage de colonnes sans
          réécriture des fichiers existants.
    \item \textbf{Time travel} : consultation de l'état des données à un
          instant passé via \texttt{AS OF}.
    \item \textbf{MERGE INTO} : opération upsert atomique combinant INSERT,
          UPDATE et DELETE en une seule instruction.
    \item \textbf{Partitionnement caché} : les requêtes n'ont pas besoin de
          connaître la stratégie de partitionnement.
\end{itemize}

\section{Debezium --- Connecteur CDC}

Debezium est un connecteur Kafka Connect open-source qui lit le WAL PostgreSQL
et publie chaque modification sur un topic Kafka au format JSON structuré.
Chaque message contient l'état \texttt{before} et \texttt{after} de la ligne,
le type d'opération et les métadonnées de la source.

\begin{lstlisting}[language={},
    caption={Structure d'un message Debezium (opération UPDATE)}]
{
  "op":     "u",
  "before": { "id": 42, "status": "pending",   "channel": "web" },
  "after":  { "id": 42, "status": "confirmed",  "channel": "web" },
  "source": {
    "db":    "ecommerce",
    "table": "orders",
    "ts_ms": 1716604800000,
    "lsn":   12589
  },
  "ts_ms": 1716604800050
}
\end{lstlisting}

\section{Stack technologique retenu}

\begin{table}[h]
\centering
\caption{Récapitulatif des technologies utilisées}
\label{tab:stack}
\renewcommand{\arraystretch}{1.3}
\begin{tabular}{lll}
\toprule
\textbf{Composant}     & \textbf{Technologie}         & \textbf{Version} \\
\midrule
Base source (OLTP)     & PostgreSQL                   & 15               \\
Bus de messages        & Apache Kafka (Confluent)     & 7.5.0            \\
Connecteur CDC         & Debezium                     & 2.5              \\
Traitement de flux     & Apache Spark Structured Streaming & 3.5.3       \\
Format de table        & Apache Iceberg               & 1.5.2            \\
Métastore              & Hive Standalone + MariaDB    & 3.1.3 + 10.11   \\
Stockage objet         & MinIO (S3-compatible)        & 2024-01-16       \\
Moteur SQL             & Trino                        & 435              \\
API REST               & FastAPI + Uvicorn            & 0.110+           \\
Authentification       & JWT (HS256) via PyJWT        & 2.8+             \\
Interface web          & React + TypeScript + Vite    & 18 + 5           \\
Visualisation BI       & Apache Superset              & 3.1.0            \\
Monitoring métriques   & Prometheus                   & 2.50.1           \\
Dashboards             & Grafana                      & 10.3.3           \\
Agrégation de logs     & Loki + Promtail              & 2.9.4            \\
Orchestration          & Docker Compose               & v2               \\
\bottomrule
\end{tabular}
\end{table}

%% ════════════════════════════════════════════════════════════
\chapter{Architecture Globale du Système}
%% ════════════════════════════════════════════════════════════

\section{Vue d'ensemble}

La figure~\ref{fig:archi} illustre l'architecture complète du pipeline CDC
Medallion. Les données transitent de gauche à droite, de la base de données
transactionnelle vers les couches analytiques.

\begin{figure}[h]
\centering
\begin{tikzpicture}[
    >=Stealth,
    box/.style={rectangle, rounded corners=5pt, minimum width=2.8cm,
                minimum height=1cm, text centered, font=\small\bfseries,
                draw, thick, inner sep=4pt},
    arr/.style={->, thick},
]
%% Colonne 1 : Source
\node[box, fill=blue!15, draw=blue!60] (pg)
    at (0,0) {PostgreSQL 15\\{\tiny WAL Logique}};

%% Colonne 2 : Ingestion
\node[box, fill=red!12, draw=red!60] (dz)
    at (3.8,1.2) {Debezium 2.5\\{\tiny Kafka Connect}};
\node[box, fill=red!20, draw=red!70] (kf)
    at (3.8,-1.2) {Apache Kafka\\{\tiny 7.5.0}};

%% Colonne 3 : Traitement
\node[box, fill=sparkblue!18, draw=sparkblue!70] (sp)
    at (7.8,0) {Apache Spark\\{\tiny 3.5.3 Streaming}};

%% Colonne 4 : Stockage
\node[box, fill=orange!18, draw=orange!60] (mn)
    at (11.5,1.2) {MinIO\\{\tiny S3 / Parquet}};
\node[box, fill=purple!15, draw=purple!60] (hv)
    at (11.5,-1.2) {Hive Metastore\\{\tiny Catalogue}};

%% Couches Iceberg (dessous)
\node[box, fill=bronze!25, draw=bronze!70] (br)
    at (3.5,-4.5) {Bronze\\{\tiny Iceberg}};
\node[box, fill=argent!25, draw=argent!70] (sv)
    at (7.8,-4.5) {Argent\\{\tiny Iceberg}};
\node[box, fill=or!35, draw=or!80] (gd)
    at (12.1,-4.5) {Or\\{\tiny Iceberg}};

%% Requête & exposition
\node[box, fill=trinoviolet!18, draw=trinoviolet!70] (tr)
    at (7.8,-7) {Trino 435\\{\tiny SQL Distribué}};
\node[box, fill=green!15, draw=green!60] (api)
    at (4,-7) {FastAPI\\{\tiny REST + JWT}};
\node[box, fill=cyan!15, draw=cyan!60] (fe)
    at (11.5,-7) {React\\{\tiny Tableau de bord}};

%% Flèches
\draw[arr, blue!60]   (pg.east) -- ++(0.4,0) |- (dz.west)
    node[pos=0.2, right, font=\tiny] {WAL};
\draw[arr, red!60]    (dz.south) -- (kf.north)
    node[midway, right, font=\tiny] {JSON};
\draw[arr, red!70]    (kf.east) -- (sp.west)
    node[midway, above, font=\tiny] {30 s batch};
\draw[arr, orange!60] (sp.east) -- ++(0.3,0) |- (mn.west)
    node[pos=0.3, above, font=\tiny] {Parquet};
\draw[arr, purple!50] (hv.west) -| (sp.south east);
\draw[arr, bronze!70] (sp.south) -- ++(0,-0.5) -| (br.north);
\draw[arr, argent!70] (sp.south) -- (sv.north);
\draw[arr, or!70]     (sp.south) -- ++(0,-0.5) -| (gd.north);
\draw[arr, trinoviolet!60] (br.south) |- (tr.west);
\draw[arr, trinoviolet!60] (sv.south) -- (tr.north);
\draw[arr, trinoviolet!60] (gd.south) |- (tr.east);
\draw[arr, green!60]  (tr.west) -- (api.east);
\draw[arr, cyan!60]   (tr.east) -- (fe.west);
\draw[arr, gray!50]   (api.north east) to[bend left=15] (fe.north west);
\end{tikzpicture}
\caption{Architecture globale du pipeline CDC Medallion (21 services)}
\label{fig:archi}
\end{figure}

\section{Flux de données end-to-end}

\begin{enumerate}
    \item Une transaction PostgreSQL est validée (INSERT / UPDATE / DELETE).
    \item Le moteur écrit l'opération dans le WAL (journal transactionnel).
    \item Debezium lit le WAL via le slot de réplication logique et publie
          un message JSON sur le topic Kafka correspondant~: latence~$<$~1~s.
    \item Spark Structured Streaming consomme les 6 topics en micro-batches
          de 30~secondes et écrit en mode \emph{append} dans la couche Bronze.
    \item Un second job Spark lit la couche Bronze et applique un
          \texttt{MERGE INTO} Iceberg pour mettre à jour la couche Argent
          (micro-batch 60~s).
    \item Un troisième job Spark agrège la couche Argent en 5 tables Or
          (cycle de 5~minutes, réécriture de la partition du jour).
    \item Trino expose les trois couches en SQL standard via le catalogue
          \texttt{iceberg}, schéma \texttt{analytics}.
    \item L'API FastAPI et le frontend React interrogent Trino pour
          les KPI, les métriques et le statut des services.
\end{enumerate}

\section{Inventaire des services Docker}

\begin{longtable}{llll}
\caption{Inventaire des 21 services Docker Compose}
\label{tab:services}\\
\toprule
\textbf{Service} & \textbf{Image} & \textbf{Version} & \textbf{Rôle} \\
\midrule
\endfirsthead
\multicolumn{4}{c}{\tablename\ \thetable{} --- suite}\\
\toprule
\textbf{Service} & \textbf{Image} & \textbf{Version} & \textbf{Rôle} \\
\midrule
\endhead
\midrule
\multicolumn{4}{r}{\itshape Suite page suivante}\\
\endfoot
\bottomrule
\endlastfoot
PostgreSQL      & postgres                                & 15       & SGBD source            \\
Zookeeper       & confluentinc/cp-zookeeper               & 7.5.0    & Coordination Kafka     \\
Kafka           & confluentinc/cp-kafka                   & 7.5.0    & Bus de messages        \\
Kafka Connect   & debezium/connect                        & 2.5      & Connecteur CDC         \\
MinIO           & minio/minio                             & 2024-01  & Stockage objet S3      \\
MinIO Init      & minio/mc (one-shot)                     & latest   & Création bucket warehouse \\
MariaDB         & mariadb                                 & 10.11    & Backend Hive           \\
Hive Metastore  & build ./hive                            & 3.1.3    & Catalogue de tables    \\
Spark Master    & apache/spark                            & 3.5.3    & Orchestrateur Spark    \\
Spark Worker    & apache/spark                            & 3.5.3    & Exécuteur Spark        \\
Trino           & trinodb/trino                           & 435      & SQL distribué          \\
Apache Superset & apache/superset                         & 3.1.0    & Visualisation BI       \\
Prometheus      & prom/prometheus                         & 2.50.1   & Collecte métriques     \\
Alertmanager    & prom/alertmanager                       & 0.27.0   & Gestion alertes        \\
Grafana         & grafana/grafana                         & 10.3.3   & Dashboards             \\
Loki            & grafana/loki                            & 2.9.4    & Indexation logs        \\
Promtail        & grafana/promtail                        & 2.9.4    & Collecte logs          \\
Kafka Exporter  & danielqsj/kafka-exporter                & latest   & Métriques Kafka        \\
Postgres Exporter & prometheuscommunity/postgres-exporter & 0.15.0   & Métriques PostgreSQL   \\
Backend API     & build ./api (Python 3.11)               & FastAPI  & API REST               \\
Frontend        & node                                    & 20-alpine & Interface React       \\
\end{longtable}

Tous les services partagent le réseau Docker bridge \texttt{cdc-net}.
Les données persistantes sont stockées dans 9 volumes Docker nommés :
\texttt{postgres-data}, \texttt{minio-data}, \texttt{mariadb-data},
\texttt{kafka-data}, \texttt{zookeeper-data}, \texttt{zookeeper-log},
\texttt{ivy2-cache}, \texttt{prometheus-data} et \texttt{grafana-data}.

%% ════════════════════════════════════════════════════════════
\chapter{Implémentation du Pipeline de Données}
%% ════════════════════════════════════════════════════════════

\section{Schéma de la base de données source}

La base PostgreSQL \texttt{ecommerce} modélise un système de commerce en
ligne avec 6 tables interconnectées.

\begin{table}[h]
\centering
\caption{Description des tables de la base source \texttt{ecommerce}}
\label{tab:schema}
\renewcommand{\arraystretch}{1.3}
\begin{tabularx}{\linewidth}{lXl}
\toprule
\textbf{Table} & \textbf{Description} & \textbf{Colonnes clés} \\
\midrule
\texttt{customers}    & Profils clients (6 colonnes)
    & id, email, city, segment \\
\texttt{products}     & Catalogue produits (6 colonnes)
    & id, category, brand, price \\
\texttt{orders}       & Commandes passées (5 colonnes)
    & id, customer\_id, status, channel \\
\texttt{order\_items} & Lignes de commande (5 colonnes)
    & order\_id, product\_id, quantity, unit\_price \\
\texttt{payments}     & Transactions de paiement (6 colonnes)
    & order\_id, method, amount, status \\
\texttt{returns}      & Retours produits (5 colonnes)
    & order\_id, reason, status \\
\bottomrule
\end{tabularx}
\end{table}

Les champs énumérés incluent :
\begin{itemize}
    \item \texttt{segment} : B2B, B2C, Premium, Standard
    \item \texttt{status} (commandes) : pending, confirmed, shipped, delivered,
          cancelled, returned
    \item \texttt{channel} : web, mobile, in-store, marketplace
    \item \texttt{method} (paiements) : credit\_card, debit\_card,
          cash\_on\_delivery, bank\_transfer, wallet
\end{itemize}

\section{Configuration de la réplication PostgreSQL}

\begin{lstlisting}[language={},
    caption={Paramètres WAL injectés via docker-compose.yml (commande postgres)}]
# Les parametres WAL sont passes en argument de demarrage du conteneur,
# ce qui evite tout ALTER SYSTEM necessitant un redemarrage.
postgres
  -c wal_level=logical
  -c max_wal_senders=4
  -c max_replication_slots=4
  -c track_commit_timestamp=on
\end{lstlisting}

\begin{lstlisting}[language=SQL,
    caption={Utilisateur Debezium et publication logique (sql/init.sql)}]
-- Utilisateur Debezium avec privilege de replication
DO $$
BEGIN
    IF NOT EXISTS (SELECT FROM pg_roles WHERE rolname = 'debezium') THEN
        CREATE USER debezium WITH PASSWORD 'debezium' REPLICATION LOGIN;
    END IF;
END
$$;

GRANT USAGE  ON SCHEMA public TO debezium;
GRANT SELECT ON ALL TABLES IN SCHEMA public TO debezium;
ALTER DEFAULT PRIVILEGES IN SCHEMA public
    GRANT SELECT ON TABLES TO debezium;

-- Publication sur les 6 tables surveillees uniquement (portee limitee)
CREATE PUBLICATION debezium_pub
    FOR TABLE customers, products, orders,
              order_items, payments, returns;
-- Note : Debezium cree et gere automatiquement son slot de replication.
\end{lstlisting}

\section{Connecteur Debezium}

Le connecteur est enregistré via l'API REST de Kafka Connect~:

\begin{lstlisting}[language={},
    caption={Configuration du connecteur Debezium (config/debezium-connector.json)}]
{
  "name": "cdc-postgres",
  "config": {
    "connector.class":
        "io.debezium.connector.postgresql.PostgresConnector",
    "database.hostname":       "postgres",
    "database.port":           "5432",
    "database.user":           "debezium",
    "database.dbname":         "ecommerce",
    "database.server.name":    "cdc",
    "plugin.name":             "pgoutput",
    "publication.name":        "debezium_pub",
    "snapshot.mode":           "initial",
    "decimal.handling.mode":   "double",
    "time.precision.mode":     "connect",
    "tombstones.on.delete":    "false",
    "heartbeat.interval.ms":   "10000",
    "key.converter.schemas.enable":   "false",
    "value.converter.schemas.enable": "false"
  }
}
\end{lstlisting}

Le connecteur produit 6 topics Kafka :
\texttt{cdc.public.customers}, \texttt{cdc.public.products},
\texttt{cdc.public.orders}, \texttt{cdc.public.order\_items},
\texttt{cdc.public.payments}, \texttt{cdc.public.returns}.

\section{Couche Bronze --- Ingestion brute}

Le job Bronze (\texttt{spark\_jobs/bronze\_job.py}, 179~lignes) consomme les
6~topics Kafka et écrit dans des tables Iceberg en mode \textbf{append-only}.
Chaque événement est conservé tel quel, sans transformation, garantissant
l'auditabilité complète.

\begin{table}[h]
\centering
\caption{Schéma des tables Bronze (colonnes principales)}
\label{tab:bronze}
\renewcommand{\arraystretch}{1.2}
\begin{tabular}{lll}
\toprule
\textbf{Colonne} & \textbf{Type Spark} & \textbf{Description} \\
\midrule
\texttt{kafka\_offset}       & BIGINT    & Offset du message Kafka \\
\texttt{kafka\_topic}        & STRING    & Nom du topic source \\
\texttt{kafka\_timestamp}    & TIMESTAMP & Heure de publication Kafka \\
\texttt{op}                  & STRING    & Opération CDC : c / u / d / r \\
\texttt{before\_payload}     & STRING    & JSON de l'état avant \\
\texttt{after\_payload}      & STRING    & JSON de l'état après \\
\texttt{source\_ts\_ms}      & BIGINT    & Horodatage WAL PostgreSQL \\
\texttt{raw\_payload}        & STRING    & Message Debezium complet \\
\texttt{ingestion\_timestamp}& TIMESTAMP & Heure d'ingestion Spark \\
\texttt{ingestion\_date}     & DATE      & \textbf{Colonne de partition} \\
\bottomrule
\end{tabular}
\end{table}

\begin{lstlisting}[language=Python,
    caption={Extrait du job Bronze --- 6 streams independants, un par table}]
TABLES = ["customers","products","orders",
          "order_items","payments","returns"]

def parse_and_write(spark, table, df, batch_id):
    parsed = (df
        .withColumn("raw_str", col("value").cast("string"))
        # schemas.enable=false : le message EST le payload (pas de wrapper).
        # Repli sur $.payload si le connecteur repasse en schemas active.
        .withColumn("payload", coalesce(
            get_json_object(col("raw_str"), "$.payload"),
            col("raw_str")))
        .withColumn("op",           get_json_object(col("payload"), "$.op"))
        .withColumn("before_payload",get_json_object(col("payload"),"$.before"))
        .withColumn("after_payload", get_json_object(col("payload"),"$.after"))
        .withColumn("source_ts_ms",
            get_json_object(col("payload"),"$.source.ts_ms").cast("long"))
        .withColumn("ingestion_timestamp", current_timestamp())
        .withColumn("ingestion_date",      current_date()))

    # Separation lignes valides / invalides en un seul passage
    # (heartbeats Debezium -> op=null, ne doivent pas etre persistes)
    valid   = parsed.filter(col("op").isNotNull())
    invalid = parsed.filter(col("op").isNull())

    dropped = invalid.count()
    if dropped > 0:
        logger.warning("Batch %d | %s | %d lignes null-op ignorees",
                       batch_id, table, dropped)

    valid.writeTo(f"iceberg.analytics.bronze_{table}_cdc").append()

def start_table_stream(spark, table):
    # Chaque table a son propre stream et son propre checkpoint.
    # Cela permet une recuperation individuelle en cas de defaillance.
    return (spark.readStream
        .format("kafka")
        .option("kafka.bootstrap.servers", "kafka:9092")
        .option("subscribe", f"cdc.public.{table}")
        .option("startingOffsets", "earliest")
        .option("maxOffsetsPerTrigger", 10_000)
        .option("failOnDataLoss", "false")
        .load()
        .writeStream
        .foreachBatch(lambda df, bid: parse_and_write(spark, table, df, bid))
        .option("checkpointLocation",
                f"s3a://warehouse/checkpoints/bronze/{table}")
        .trigger(processingTime="30 seconds")
        .start())

# Lancement de 6 streams independants, un par table source
[start_table_stream(spark, t) for t in TABLES]
spark.streams.awaitAnyTermination()
\end{lstlisting}

\section{Couche Argent --- Déduplication et typage fort}

Le job Argent (\texttt{spark\_jobs/silver\_job.py}, 252~lignes) lit en
continu les tables Bronze et applique un \texttt{MERGE INTO} Iceberg pour
maintenir l'état courant de chaque entité (micro-batch de 60~secondes).

\subsection{Déduplication par fenêtre}

\begin{lstlisting}[language=Python,
    caption={Déduplication --- conservation de la dernière version par PK}]
from pyspark.sql.window import Window
from pyspark.sql.functions import row_number, desc

# Pour chaque clé primaire, conserver uniquement
# l'événement le plus récent (par source_ts_ms)
window_spec = Window.partitionBy("id").orderBy(desc("_source_ts_ms"))

df_dedup = (df_parsed
    .withColumn("_rn", row_number().over(window_spec))
    .filter(col("_rn") == 1)
    .drop("_rn"))
\end{lstlisting}

\subsection{Opération MERGE INTO Iceberg}

\begin{lstlisting}[language=Python,
    caption={Upsert atomique Iceberg via MERGE INTO (vue global\_temp)}]
# La vue est enregistree en global_temp (et non en session temp).
# Raison : spark.sql.defaultCatalog=iceberg fait que Spark cherche
# les vues non qualifiees dans le catalogue Iceberg, ce qui echoue.
# global_temp.{view} contourne ce probleme de resolution de catalogue.
view = f"_silver_src_{table}_{batch_id}"
deduped.select(all_cols).createOrReplaceGlobalTempView(view)

spark.sql(f"""
    MERGE INTO iceberg.analytics.silver_{table} t
    USING global_temp.{view} s
    ON    t.{pk} = s.{pk}
    WHEN MATCHED AND s._is_deleted = true THEN
        UPDATE SET t._is_deleted = true,
                   t._op         = 'd',
                   t._updated_at = s._updated_at,
                   t.updated_date = s.updated_date
    WHEN MATCHED THEN
        UPDATE SET {set_sql}
    WHEN NOT MATCHED AND s._is_deleted = false THEN
        INSERT ({ins_cols}) VALUES ({ins_vals})
""")
\end{lstlisting}

Les tables Argent ajoutent 4 colonnes de métadonnées à chaque ligne :
\texttt{\_op}, \texttt{\_source\_ts\_ms}, \texttt{\_updated\_at},
\texttt{\_is\_deleted}, ainsi qu'une colonne de partition \texttt{updated\_date}.

\section{Couche Or --- Agrégats métier}

Le job Or (\texttt{spark\_jobs/gold\_job.py}, 364~lignes) produit
5~tables analytiques en agrégeant la couche Argent toutes les 5~minutes.

\begin{table}[h]
\centering
\caption{Tables de la couche Or et leurs indicateurs}
\label{tab:gold}
\renewcommand{\arraystretch}{1.3}
\begin{tabularx}{\linewidth}{lXl}
\toprule
\textbf{Table} & \textbf{Indicateurs produits} & \textbf{Granularité} \\
\midrule
\texttt{gold\_customers}
    & Nombre de commandes, dépenses totales, première/dernière commande
    & Par client \\
\texttt{gold\_orders}
    & Commande enrichie : client, paiement, retour, total
    & Par commande \\
\texttt{gold\_sales\_facts}
    & Revenus bruts/nets, volumes vendus, taux d'annulation/retour
    & Date $\times$ canal $\times$ segment $\times$ catégorie \\
\texttt{gold\_payments\_summary}
    & Nombre et montant par méthode de paiement
    & Mois $\times$ méthode \\
\texttt{gold\_returns\_summary}
    & Nombre et taux de retour par raison
    & Mois $\times$ raison \\
\bottomrule
\end{tabularx}
\end{table}

\begin{lstlisting}[language=Python,
    caption={Agrégation des faits de vente --- gold\_sales\_facts}]
gold_sales = (silver_order_items
    .join(silver_products.select(
              col("id").alias("prid"), "category"),
          silver_order_items.product_id == col("prid"), "left")
    .join(silver_orders.select(
              col("id").alias("oid"),
              "order_date", "customer_id", "status", "channel"),
          silver_order_items.order_id == col("oid"), "left")
    .join(silver_customers.select(
              col("id").alias("cid"), "segment"),
          col("customer_id") == col("cid"), "left")
    .groupBy(
        col("order_date").alias("sale_date"), "channel",
        col("segment").alias("customer_segment"),
        col("category").alias("product_category"))
    .agg(
        countDistinct(silver_order_items.order_id).alias("order_count"),
        # countDistinct par statut : evite de compter plusieurs fois
        # la meme commande qui a plusieurs lignes d'articles.
        countDistinct(when(col("status") == "delivered",
            silver_order_items.order_id)).alias("delivered_count"),
        countDistinct(when(col("status") == "cancelled",
            silver_order_items.order_id)).alias("cancelled_count"),
        countDistinct(when(col("status") == "returned",
            silver_order_items.order_id)).alias("returned_count"),
        F.sum(col("quantity") * col("unit_price")).alias("gross_revenue"),
        F.sum(when(col("status").isin("delivered","shipped"),
              col("quantity") * col("unit_price")).otherwise(0.0)
             ).alias("net_revenue"),
        F.sum("quantity").alias("total_items_sold"),
    )
    # avg_order_value = revenu_brut / nb_commandes (pas avg des lignes)
    .withColumn("avg_order_value",
        when(col("order_count") > 0,
             col("gross_revenue") / col("order_count")).otherwise(0.0))
    .withColumn("_updated_at",   current_timestamp())
    .withColumn("snapshot_date", current_date()))

# Réécriture de la partition du jour (idempotent, cycle 5 min)
gold_sales.writeTo("iceberg.analytics.gold_sales_facts").overwritePartitions()
\end{lstlisting}

\section{Configuration Iceberg et connexion S3A/MinIO}

\begin{lstlisting}[caption={spark-defaults.conf --- Catalogue Iceberg sur Hive + MinIO}]
# Extensions Iceberg pour Spark SQL
spark.sql.extensions=\
    org.apache.iceberg.spark.extensions.IcebergSparkSessionExtensions

# Catalogue Iceberg (type Hive, métastore MariaDB)
spark.sql.catalog.iceberg=org.apache.iceberg.spark.SparkCatalog
spark.sql.catalog.iceberg.type=hive
spark.sql.catalog.iceberg.uri=thrift://hive-metastore:9083
spark.sql.catalog.iceberg.warehouse=s3a://warehouse/
spark.sql.defaultCatalog=iceberg

# Connexion MinIO via protocole S3A
spark.hadoop.fs.s3a.endpoint=http://minio:9000
spark.hadoop.fs.s3a.path.style.access=true
spark.hadoop.fs.s3a.impl=org.apache.hadoop.fs.s3a.S3AFileSystem
spark.hadoop.fs.s3a.connection.ssl.enabled=false
\end{lstlisting}

\section{Requêtes KPI via Trino}

Trino expose les trois couches sous le catalogue \texttt{iceberg},
schéma \texttt{analytics}. Exemples de requêtes analytiques :

\begin{lstlisting}[language=SQL,
    caption={KPI analytiques exécutés via Trino 435}]
-- Revenus par canal (14 derniers jours)
SELECT   channel,
         SUM(gross_revenue) AS revenu_brut,
         SUM(net_revenue)   AS revenu_net
FROM     iceberg.analytics.gold_sales_facts
WHERE    sale_date >= current_date - INTERVAL '14' DAY
GROUP BY channel
ORDER BY revenu_brut DESC;

-- Mix des méthodes de paiement (mois courant)
SELECT   method AS methode,
         SUM(transaction_count) AS nb_transactions,
         ROUND(SUM(total_amount), 2) AS montant_total
FROM     iceberg.analytics.gold_payments_summary
WHERE    payment_month = FORMAT_DATETIME(NOW(), 'yyyy-MM')
GROUP BY method;

-- Latence CDC end-to-end (percentile 95) -- couche Bronze
SELECT   source_table,
         ROUND(
           approx_percentile(
             (to_unixtime(ingestion_timestamp)*1000 - source_ts_ms)/1000.0,
             0.95),
         2) AS latence_p95_secondes
FROM     iceberg.analytics.bronze_orders_cdc
GROUP BY source_table;
\end{lstlisting}

%% ════════════════════════════════════════════════════════════
\chapter{API REST, Interface Utilisateur et Monitoring}
%% ════════════════════════════════════════════════════════════

\section{API REST --- FastAPI}

L'API REST est développée avec \textbf{FastAPI~0.110+} et Uvicorn,
en Python~3.11, exposée sur le port 8000.

\subsection{Authentification par token JWT}

À la connexion, le serveur émet un token JWT signé (algorithme HS256,
validité 60~minutes). Ce token est transmis dans le header
\texttt{Authorization:~Bearer~<token>} pour toutes les requêtes protégées.

\subsection{Endpoints}

\begin{table}[h]
\centering
\caption{Endpoints de l'API REST (13 endpoints, 7 routeurs)}
\label{tab:api}
\renewcommand{\arraystretch}{1.3}
\begin{tabularx}{\linewidth}{llXl}
\toprule
\textbf{Méthode} & \textbf{Endpoint} & \textbf{Description} & \textbf{Auth} \\
\midrule
POST  & \texttt{/auth/login}
    & Connexion, retourne un JWT & Non \\
POST  & \texttt{/auth/logout}
    & Déconnexion (stateless, côté client) & Non \\
GET   & \texttt{/auth/me}
    & Profil de l'utilisateur connecté & Oui \\
\midrule
GET   & \texttt{/api/health}
    & Statut et latence des 5 services & Non \\
GET   & \texttt{/api/metrics/\{layer\}}
    & Comptage de lignes par table Iceberg (cache 30~s) & Oui \\
POST  & \texttt{/api/query}
    & SQL SELECT vers Trino, réponse NDJSON streamée & Oui \\
GET   & \texttt{/api/events/stream}
    & Flux SSE temps réel des événements Kafka CDC & Oui \\
\midrule
GET   & \texttt{/api/connector/status}
    & Statut du connecteur Debezium & Non \\
POST  & \texttt{/api/connector/restart}
    & Redémarrage du connecteur Debezium & Admin \\
GET   & \texttt{/api/spark/jobs}
    & 20 derniers jobs (Spark History Server) & Non \\
GET   & \texttt{/api/spark/running}
    & Applications actives sur le Spark Master & Non \\
\midrule
GET   & \texttt{/api/alerts}
    & Liste des alertes pipeline & Oui \\
PATCH & \texttt{/api/alerts/\{id\}/ack}
    & Acquittement d'alerte & Oui \\
\bottomrule
\end{tabularx}
\end{table}

\subsection{Sécurisation des requêtes SQL}

L'endpoint \texttt{/api/query} valide le SQL avant transmission à Trino :
toute instruction de modification (\texttt{DROP}, \texttt{CREATE},
\texttt{INSERT}, \texttt{UPDATE}, \texttt{DELETE}, \texttt{ALTER},
\texttt{TRUNCATE}, \texttt{GRANT}, \texttt{MERGE}) déclenche une erreur
HTTP~400. Seuls les \texttt{SELECT} sont autorisés.

\subsection{Streaming NDJSON}

Les résultats volumineux sont renvoyés en \textit{Newline-Delimited JSON}
via \texttt{StreamingResponse}, évitant de charger toute la réponse en mémoire.

\subsection{Flux SSE --- Kafka en temps réel}

L'endpoint \texttt{GET /api/events/stream} ouvre un consommateur Kafka
\textit{par connexion} avec un \texttt{group\_id} unique
(\texttt{cdc-ui-consumer-<uuid>}). Ce choix garantit que chaque onglet
navigateur reçoit l'intégralité des événements, sans que Kafka les
répartisse entre clients. Un heartbeat \texttt{: keep-alive} est émis
toutes les 15~secondes pour maintenir la connexion ouverte.

\subsection{Contrainte d'unicité du worker Uvicorn}

Le Dockerfile de production lance un \textbf{unique} worker Uvicorn
(sans l'option \texttt{--workers}). Ce choix est délibéré : plusieurs
workers maintiendraient chacun un cache mémoire indépendant pour l'endpoint
\texttt{/api/metrics}, multipliant les requêtes Trino inutiles. Pour
passer à l'échelle, il convient de déployer plusieurs conteneurs derrière
un équilibreur de charge et de remplacer le cache en mémoire par Redis.

\begin{lstlisting}[language=Python,
    caption={Endpoint /api/query --- streaming NDJSON}]
@router.post("/query")
async def run_query(body: QueryRequest,
                    user: User = Depends(get_current_user)):
    validate_read_only_sql(body.sql)    # Bloque DDL/DML

    async def row_generator():
        async for row in trino_client.query_stream(body.sql):
            yield json.dumps(row, default=str) + "\n"

    return StreamingResponse(row_generator(),
                             media_type="application/x-ndjson")
\end{lstlisting}

\section{Interface Utilisateur React}

Le frontend est une application monopage développée en React~18 avec
TypeScript, Vite et Tailwind CSS (thème sombre), organisée en 6 routes :

\begin{table}[h]
\centering
\caption{Pages de l'interface utilisateur (6 routes)}
\label{tab:pages}
\renewcommand{\arraystretch}{1.3}
\begin{tabularx}{\linewidth}{llXl}
\toprule
\textbf{Route} & \textbf{Page} & \textbf{Contenu} & \textbf{Rafraîchissement} \\
\midrule
\texttt{/dashboard}
    & KPI Overview
    & KPIs animés (count-up), flux Kafka simulé, statut
      des services, cartes Medallion avec jauges de volume
    & Auto 30 s \\
\texttt{/bronze}
    & Bronze Layer
    & Liste des tables Iceberg avec comptage de lignes,
      date de partition, taille ; tiroir de schéma latéral
    & Manuel \\
\texttt{/silver}
    & Silver Layer
    & Idem Bronze — données dédupliquées et typées
    & Manuel \\
\texttt{/gold}
    & Gold Layer
    & Idem Bronze — agrégats métier prêts à l'analyse
    & Manuel \\
\texttt{/query}
    & Query Console
    & Éditeur SQL Monaco (thème sombre), exécution vers Trino,
      résultats tabulaires streamés en NDJSON, export CSV
    & À la demande \\
\texttt{/stream}
    & Event Stream
    & Flux SSE en direct des événements CDC Kafka,
      filtre par table et topic, vues Tableau et Timeline
    & Temps réel (SSE) \\
\bottomrule
\end{tabularx}
\end{table}

\section{Monitoring et Observabilité}

La stack de supervision repose sur quatre composants :

\begin{itemize}
    \item \textbf{Prometheus} (port~9090) : collecte toutes les 15~secondes
          les métriques de Kafka, PostgreSQL, MinIO, Spark et de l'API.
    \item \textbf{Grafana} (port~3001) : dashboards préconfigurés pour
          visualiser les métriques Prometheus et les logs Loki.
    \item \textbf{Loki + Promtail} : agrégation et indexation des logs Docker
          de tous les services.
    \item \textbf{Alertmanager} (port~9093) : envoi d'alertes (Slack ou e-mail)
          lorsqu'un seuil est dépassé.
\end{itemize}

\begin{table}[h]
\centering
\caption{Règles d'alerte Prometheus configurées}
\label{tab:alerts}
\renewcommand{\arraystretch}{1.3}
\begin{tabular}{ll}
\toprule
\textbf{Condition} & \textbf{Seuil} \\
\midrule
Lag consommateur Kafka & $>$ 10\,000 messages en retard \\
Latence couche Bronze & $>$ 120 s \\
Disponibilité des services & $<$ 100\,\% \\
Espace disque MinIO & $>$ 80\,\% \\
Taux d'erreur API & $>$ 5\,\% sur 1 minute \\
\bottomrule
\end{tabular}
\end{table}

%% ════════════════════════════════════════════════════════════
\chapter{Déploiement et Résultats}
%% ════════════════════════════════════════════════════════════

\section{Environnement de déploiement}

\begin{table}[h]
\centering
\caption{Spécifications de la machine virtuelle Oracle Cloud}
\label{tab:vm}
\renewcommand{\arraystretch}{1.3}
\begin{tabular}{ll}
\toprule
\textbf{Paramètre}       & \textbf{Valeur} \\
\midrule
Fournisseur              & Oracle Cloud Infrastructure (OCI) \\
Système d'exploitation   & Ubuntu 22.04 LTS \\
RAM totale               & 11,67 Go \\
vCPU                     & 4 \\
Stockage                 & 50 Go (SSD) \\
Orchestration            & Docker Compose v2 (21 services) \\
Dépôt source             & GitHub (git pull \& restart) \\
\bottomrule
\end{tabular}
\end{table}

\section{Gestion de la mémoire}

La contrainte principale est la coexistence de plusieurs processus JVM
lourds sur une VM de 11,67~Go. La configuration initiale provoquait
2,6~Go de swap, multipliant les temps de réponse par 5 à 10.

\begin{table}[h]
\centering
\caption{Limites mémoire par service après optimisation}
\label{tab:memory}
\renewcommand{\arraystretch}{1.3}
\begin{tabular}{lll}
\toprule
\textbf{Service} & \textbf{Limite conteneur} & \textbf{Paramètre JVM} \\
\midrule
Trino           & 3,0 Go  & \texttt{-Xmx2G -XX:ReservedCodeCacheSize=256M} \\
Spark Worker    & 1,5 Go  & \texttt{SPARK\_WORKER\_MEMORY=1g}               \\
Kafka           & 768 Mo  & \texttt{KAFKA\_HEAP\_OPTS=-Xmx512M}             \\
Kafka Connect   & 640 Mo  & Paramètre par défaut Debezium                  \\
Hive Metastore  & 700 Mo  & Paramètre par défaut Hive                      \\
\bottomrule
\end{tabular}
\end{table}

Le problème Trino résultait d'une contrainte de configuration : la somme de
\texttt{query.max-memory-per-node} (1~Go) et de
\texttt{memory.heap-headroom-per-node} (256~Mo) doit être strictement
inférieure à la taille du heap JVM (\texttt{-Xmx2G}). En déploiement
mono-nœud, \texttt{query.max-memory-per-node} plafonne toujours en premier,
ce qui rend inutile de fixer \texttt{query.max-memory} au-dessus de 1~Go.
La configuration retenue :

\begin{lstlisting}[caption={config/trino/config.properties --- limites mémoire}]
query.max-memory=1GB
query.max-memory-per-node=1GB
memory.heap-headroom-per-node=256MB
# Deploiement mono-noeud : query.max-memory-per-node plafonne toujours
# en premier. Fixer query.max-memory > 1 GB serait sans effet et trompeur.
# Contrainte respectee : 1 Go + 256 Mo = 1,25 Go < 2 Go (Xmx JVM).
\end{lstlisting}

\section{Séquence de démarrage}

\begin{table}[h]
\centering
\caption{Phases de démarrage du pipeline (premier lancement)}
\label{tab:startup}
\renewcommand{\arraystretch}{1.3}
\begin{tabular}{lll}
\toprule
\textbf{Phase} & \textbf{Durée estimée} & \textbf{Événement} \\
\midrule
Démarrage conteneurs   & 1 -- 2 min    & Health checks Docker passés     \\
Init Hive Metastore    & 1 -- 2 min    & Schéma MariaDB créé             \\
Démarrage Trino        & 30 -- 60 s    & Moteur SQL prêt                 \\
Téléchargement JARs    & 2 -- 5 min    & Iceberg / Kafka / S3A (Maven)   \\
Snapshot Debezium      & 20 -- 30 s    & Lecture initiale de la DB       \\
Job Bronze actif       & 30 s          & Premiers micro-batches          \\
Job Argent actif       & 60 s          & Premier MERGE INTO              \\
Job Or actif           & 5 min         & Première agrégation             \\
\midrule
\textbf{Total}         & \textbf{8 -- 15 min} & \textbf{Pipeline opérationnel} \\
\bottomrule
\end{tabular}
\end{table}

\section{Résultats et métriques de performance}

\subsection{Latences mesurées}

\begin{table}[h]
\centering
\caption{Latences end-to-end après optimisation}
\label{tab:latency}
\renewcommand{\arraystretch}{1.3}
\begin{tabular}{lll}
\toprule
\textbf{Étape} & \textbf{Latence mesurée} & \textbf{Objectif initial} \\
\midrule
PostgreSQL WAL $\to$ Kafka    & $<$ 1 s        & $<$ 5 s    \\
Kafka $\to$ Bronze (Iceberg)  & 15 -- 30 s     & $<$ 60 s   \\
Bronze $\to$ Argent (MERGE)   & 30 -- 60 s     & $<$ 120 s  \\
Argent $\to$ Or (agrégation)  & 3 -- 5 min     & $<$ 10 min \\
\textbf{End-to-end total}     & \textbf{4 -- 7 min} & \textbf{$<$ 15 min} \\
\bottomrule
\end{tabular}
\end{table}

\subsection{Données traitées}

\begin{itemize}
    \item \textbf{6 tables source} avec $\approx$100 lignes initiales par table
          (chargées par \texttt{sql/seed.sql}).
    \item \textbf{600 événements Kafka} au snapshot initial, puis flux continu
          généré par le script \texttt{scripts/simulate\_changes.sh}.
    \item \textbf{Couche Bronze} : 6 tables Iceberg partitionnées par
          \texttt{ingestion\_date}.
    \item \textbf{Couche Argent} : 6 tables avec état courant dédupliqué,
          partitionnées par \texttt{updated\_date}.
    \item \textbf{Couche Or} : 5 tables agrégées prêtes pour les requêtes
          analytiques, partitionnées par \texttt{snapshot\_date}.
\end{itemize}

\subsection{Interface opérationnelle}

L'interface React déployée offre une vision temps réel du pipeline :
\begin{itemize}
    \item Statut animé (ping vert / amber) de chaque service avec latence
          en milliseconde.
    \item Comptage de lignes par table pour les trois couches avec barres
          de volume normalisées.
    \item Graphiques de KPI : revenus sur 14~jours, répartition par canal,
          mix de paiements, performance par segment client.
    \item Topologie du pipeline avec état de santé en temps réel de chaque
          nœud (PostgreSQL $\to$ Debezium $\to$ Kafka $\to$ Spark $\to$
          MinIO $\to$ Bronze $\to$ Argent $\to$ Or).
\end{itemize}

%% ════════════════════════════════════════════════════════════
\chapter{Conclusion et Perspectives}
%% ════════════════════════════════════════════════════════════

\section{Bilan}

Ce projet démontre qu'il est possible d'implémenter un pipeline CDC
professionnel --- complet, opérationnel et reproductible --- sur une
infrastructure mono-machine à faible coût. Les principaux acquis sont :

\begin{enumerate}
    \item \textbf{Maîtrise du CDC par lecture de log} : compréhension du WAL
          PostgreSQL, de la réplication logique et de son intégration avec
          Debezium et Kafka.

    \item \textbf{Architecture Medallion} : conception des trois couches avec
          compromis explicites entre latence, fiabilité et richesse sémantique.

    \item \textbf{Spark Structured Streaming} : utilisation avancée des
          micro-batches, du \texttt{MERGE INTO} Iceberg et de la gestion des
          checkpoints pour la tolérance aux pannes.

    \item \textbf{Gestion JVM en production} : interactions complexes entre
          heap JVM, CodeCache et limites de conteneurs Docker --- problème réel
          rencontré et résolu lors du déploiement.

    \item \textbf{Observabilité complète} : stack Prometheus + Grafana + Loki
          pour les métriques, dashboards et logs de 21 services.

    \item \textbf{API et interface} : exposition sécurisée des données via
          FastAPI (JWT, streaming NDJSON, SSE Kafka) et tableau de bord
          React en thème sombre avec 6 vues opérationnelles.

    \item \textbf{Fiabilité du code} : suite de 56 tests automatisés
          (pytest) couvrant les 7 routeurs de l'API avec des mocks corrects
          du gestionnaire de contexte asynchrone AsyncPG ; validation de
          sécurité des secrets au démarrage.

    \item \textbf{Robustesse du job Bronze} : suppression du pattern
          \texttt{cache()/count()/writeTo()} qui matérialisait l'intégralité
          du micro-batch en mémoire exécuteur, source potentielle d'OOM ;
          remplacement par un filtrage en un seul passage avec journalisation
          des événements écartés (heartbeats Debezium).
\end{enumerate}

\section{Difficultés rencontrées}

\begin{itemize}
    \item \textbf{Pression mémoire} : la cohabitation de plusieurs JVM lourds
          (Trino, Spark, Kafka) sur 11,67~Go a causé 2,6~Go de swap, dégradant
          fortement les performances. Solution : réduction fine des paramètres
          \texttt{-Xmx}, \texttt{ReservedCodeCacheSize} et des limites
          \texttt{query.max-memory-per-node}.

    \item \textbf{Démarrage ordonné} : les dépendances entre services (Kafka
          avant Debezium, Hive Metastore avant Spark et Trino) nécessitent des
          health checks Docker précis et des scripts de vérification.

    \item \textbf{Compatibilité des versions} : la combinaison Spark~3.5.3 +
          Iceberg~1.5.2 + Hadoop AWS~3.3.4 + AWS SDK~1.12.262 a demandé un
          alignement précis des versions pour éviter les conflits de classpath.

    \item \textbf{Téléchargement des JARs} : le premier démarrage de Spark
          télécharge $\approx$500~Mo de dépendances Maven. Un volume partagé
          \texttt{ivy2-cache} a été ajouté pour éviter ce téléchargement à
          chaque redémarrage.
\end{itemize}

\section{Perspectives d'amélioration}

\begin{enumerate}
    \item \textbf{CI/CD GitHub Actions} : pipeline lint + test + build
          d'image Docker déclenché à chaque pull request ; déploiement
          automatique sur la VM OCI à chaque merge sur \texttt{main}.

    \item \textbf{Déploiement Kubernetes} : migration vers un cluster K8s
          (manifestes Kustomize et chart Helm déjà présents dans \texttt{k8s/}
          et \texttt{helm/}) pour la scalabilité horizontale et la haute
          disponibilité. Prévoir Strimzi Kafka Operator et Spark-on-K8s.

    \item \textbf{Optimisation Iceberg v2} : utilisation des suppressions au
          niveau ligne (\textit{position deletes}) et du \textit{z-ordering}
          pour accélérer les requêtes analytiques sur les grandes tables Or.
          Planifier les jobs de compaction (\texttt{rewrite\_data\_files},
          \texttt{expire\_snapshots}) pour contenir la croissance des petits
          fichiers Parquet.

    \item \textbf{Qualité des données} : intégration d'un framework de
          validation (Great Expectations, dbt) entre les couches Argent et Or
          pour détecter les anomalies (régression du nombre de lignes $>20$\,\%,
          clés primaires nulles, latence SLA) dès leur apparition.

    \item \textbf{Dead Letter Queue} : router les événements Debezium
          malformés vers un topic \texttt{\{table\}.dlq} au lieu de les
          ignorer silencieusement, afin d'offrir une traçabilité complète des
          données écartées.

    \item \textbf{Machine Learning} : exposition de la couche Or à MLflow pour
          construire des modèles prédictifs (churn client, détection de fraude,
          recommandation de produits).

    \item \textbf{Multi-sources} : extension à d'autres bases (MySQL, MongoDB)
          via des connecteurs Debezium supplémentaires.

    \item \textbf{Cache distribué} : remplacement du cache mémoire de
          l'endpoint \texttt{/api/metrics} par Redis, permettant le passage à
          plusieurs workers Uvicorn ou conteneurs API en parallèle.
\end{enumerate}

%% ════════════════════════════════════════════════════════════
%% Bibliographie
%% ════════════════════════════════════════════════════════════
\begin{thebibliography}{99}

\bibitem{debezium}
Debezium Authors,
\textit{Debezium Documentation --- PostgreSQL Connector},
\url{https://debezium.io/documentation/reference/stable/connectors/postgresql.html},
2024.

\bibitem{iceberg}
Apache Iceberg PMC,
\textit{Apache Iceberg --- The Open Table Format for Analytic Datasets},
\url{https://iceberg.apache.org/docs/latest/},
2024.

\bibitem{spark}
Apache Spark PMC,
\textit{Structured Streaming Programming Guide},
\url{https://spark.apache.org/docs/3.5.3/structured-streaming-programming-guide.html},
2024.

\bibitem{trino}
Trino Foundation,
\textit{Trino --- Distributed SQL Query Engine},
\url{https://trino.io/docs/current/},
2024.

\bibitem{kafka}
Apache Kafka PMC,
\textit{Apache Kafka Documentation},
\url{https://kafka.apache.org/documentation/},
2024.

\bibitem{medallion}
Databricks,
\textit{What is the Medallion Lakehouse Architecture?},
\url{https://www.databricks.com/glossary/medallion-architecture},
2024.

\bibitem{fastapi}
S. Ramírez,
\textit{FastAPI --- Modern, Fast Web Framework for Building APIs with Python},
\url{https://fastapi.tiangolo.com/},
2024.

\bibitem{prometheus}
Prometheus Authors,
\textit{Prometheus --- Monitoring System and Time Series Database},
\url{https://prometheus.io/docs/introduction/overview/},
2024.

\bibitem{kleppmann}
M. Kleppmann,
\textit{Designing Data-Intensive Applications},
O'Reilly Media, 2017, ch.~11, pp.~453--490.

\bibitem{minio}
MinIO Inc.,
\textit{MinIO Object Storage Documentation},
\url{https://min.io/docs/minio/linux/index.html},
2024.

\end{thebibliography}

%% ════════════════════════════════════════════════════════════
%% Annexes
%% ════════════════════════════════════════════════════════════
\appendix

\chapter{Commandes de démarrage}

\begin{lstlisting}[language=bash,
    caption={Démarrage complet du pipeline CDC Medallion}]
# 1. Cloner le dépôt
git clone https://github.com/Elghouali/cdc-pipeline-bigdata.git
cd cdc-pipeline-bigdata

# 2. Configurer l'environnement
cp .env.example .env
# Éditer .env : mots de passe, secrets JWT, etc.

# 3. Démarrer tous les services Docker (20 conteneurs)
docker compose up -d

# 4. Vérifier l'état des services
docker compose ps

# 5. Enregistrer le connecteur Debezium (après ~60 s)
bash scripts/register_connector.sh

# 6. Lancer les trois jobs Spark Structured Streaming
for job in bronze_job silver_job gold_job; do
    docker exec -d cdc-spark-master \
        /opt/spark/bin/spark-submit \
        --master spark://spark-master:7077 \
        --conf spark.cores.max=2 \
        /opt/spark-apps/${job}.py
done

# 7. Simuler des modifications de données en continu
bash scripts/simulate_changes.sh

# 8. Accéder à l'interface web (http://<IP>:3000)
# Identifiants : admin / <valeur API_ADMIN_PASSWORD dans .env>
\end{lstlisting}

\chapter{Arborescence du projet}

\begin{lstlisting}[caption={Structure du dépôt cdc-pipeline-bigdata}]
cdc-pipeline/
├── api/                        # API REST FastAPI (Python 3.11)
│   ├── main.py                 # Point d'entrée, lifespan, CORS
│   ├── auth.py                 # JWT : émission et vérification (PyJWT)
│   ├── config.py               # Settings + validation sécurité au démarrage
│   ├── trino_client.py         # Client Trino asynchrone (httpx)
│   ├── db.py                   # Pool AsyncPG, users, alerts
│   ├── routers/                # 7 modules d'endpoints (13 routes)
│   │   ├── health.py           # GET /api/health
│   │   ├── query.py            # POST /api/query (NDJSON)
│   │   ├── metrics.py          # GET /api/metrics/{layer} (cache 30s)
│   │   ├── alerts.py           # GET/PATCH /api/alerts
│   │   ├── events.py           # GET /api/events/stream (SSE Kafka)
│   │   ├── connector.py        # GET|POST /api/connector
│   │   └── spark.py            # GET /api/spark/jobs|running
│   ├── tests/                  # Suite de tests (56 tests, pytest)
│   │   ├── test_health.py      # Santé + root endpoint
│   │   ├── test_auth.py        # Login, logout, /auth/me
│   │   ├── test_metrics.py     # Métriques par layer + cache
│   │   ├── test_query.py       # Validation SQL + streaming
│   │   ├── test_connector.py   # Statut + redémarrage connecteur
│   │   ├── test_alerts.py      # Liste + acquittement alertes
│   │   └── test_spark.py       # Jobs + running Spark
│   └── Dockerfile              # python:3.11-slim, 1 worker Uvicorn
├── config/
│   ├── hive/hive-site.xml      # Metastore + S3A
│   ├── spark/spark-defaults.conf  # Iceberg + S3A + JARs
│   └── trino/
│       ├── config.properties   # 1 Go/noeud, port 8080, discovery
│       ├── jvm.config          # -Xmx2G, GC, CodeCache
│       └── catalog/iceberg.properties
├── frontend/                   # React 18 + TypeScript + Tailwind (thème sombre)
│   └── src/
│       ├── pages/              # 6 routes : dashboard, layers, query, stream
│       ├── components/         # Sidebar, TopBar, Badge, Button...
│       └── index.css           # Design system sombre, palette Claude
├── monitoring/                 # Prometheus, Grafana, Loki, Alertmanager
├── spark_jobs/
│   ├── bronze_job.py           # 179 lignes : Kafka → Bronze (6 streams)
│   ├── silver_job.py           # 252 lignes : Bronze → Argent (MERGE INTO)
│   └── gold_job.py             # 364 lignes : Argent → Or (5 tables)
├── sql/
│   ├── init.sql                # DDL, réplication, Debezium user, publication
│   └── seed.sql                # Données initiales (~100 lignes par table)
├── scripts/                    # Automatisation opérationnelle
│   ├── register_connector.sh   # Enregistrement connecteur Debezium
│   ├── simulate_changes.sh     # Génération de flux CDC continu
│   └── start.sh                # Démarrage ordonné complet
├── k8s/                        # Manifestes Kubernetes (Kustomize)
├── helm/                       # Chart Helm cdc-pipeline
└── docker-compose.yml          # 21 services, réseau cdc-net, 9 volumes
\end{lstlisting}

\chapter{Glossaire}

\begin{description}[style=nextline, leftmargin=4cm]
    \item[CDC]
        \textit{Change Data Capture} --- technique de capture des
        modifications d'une base relationnelle en temps réel.

    \item[WAL]
        \textit{Write-Ahead Log} --- journal transactionnel de PostgreSQL
        utilisé pour la réplication logique.

    \item[Iceberg]
        Format de table analytique ouvert garantissant les propriétés ACID,
        l'évolution de schéma et le time-travel.

    \item[Lakehouse]
        Architecture combinant la flexibilité d'un \textit{data lake}
        (stockage objet bas coût) avec les garanties ACID d'un entrepôt.

    \item[Medallion]
        Architecture en trois couches (Bronze / Argent / Or) de qualité
        croissante.

    \item[Micro-batch]
        Modèle de traitement de flux où les données sont traitées par
        petits intervalles de temps réguliers (30 s, 60 s...).

    \item[MERGE INTO]
        Instruction SQL combinant INSERT, UPDATE et DELETE en une seule
        opération atomique --- clé de la couche Argent.

    \item[NDJSON]
        \textit{Newline-Delimited JSON} --- format de streaming de données
        JSON, une ligne par objet.

    \item[JWT]
        \textit{JSON Web Token} --- standard d'authentification par token
        signé numériquement.

    \item[KPI]
        \textit{Key Performance Indicator} --- indicateur clé de performance
        métier.

    \item[Trino]
        Moteur de requête SQL distribué, successeur de Presto, capable
        d'interroger plusieurs sources de données hétérogènes.

    \item[S3A]
        Implémentation Hadoop du protocole Amazon S3 permettant à Spark
        d'accéder à MinIO comme à un stockage objet natif.

    \item[pgoutput]
        Plugin de sortie logique natif PostgreSQL, utilisé par Debezium
        pour décoder les changements WAL sans extension tierce.
\end{description}

\end{document}
